\section{Wissen \& Konzepte}
\slide{Definitionen von Wissen \& Konzepten}{
    \begin{itemize}
        \item \b{Wissen (Philosophie):} Wahre, gerechtfertigte Meinung (justified true belief); Unterscheidung zwischen Wissen (Fakt ist wahr) und Glauben (subjektive Überzeugung).
        \item \b{Wissen (KW):} Mentales Modell/Repräsentation, das Welt- und Selbstmodell konstituiert; partiell und dient der effizienten Informationsverarbeitung.
        \item \b{Konzept:} Elementare Einheit von Wissen; abstrakte mentale Repräsentation von Objekten (entspricht oft Wörtern/Sprache).
        \item \b{Semiotisches Dreieck:} Beziehung zwischen Objekt (Reale Welt), Konzept (Gedanke) und Zeichen (Wort).
    \end{itemize}
}
\slide{Repräsentations-Formate}{
    \begin{itemize}
        \item \b{Propositionale Repräsentation:} Symbolische, sprachähnliche Formalisierung (Logik-kompatibel); arbiträre (willkürliche) Relation zwischen Symbol und Objekt.
        \item \b{Analoge Repräsentation:} Nicht-symbolische, bildhafte Vorstellung; direkte Darstellung (z. B. räumliche Lage), oft wahrnehmungsnah.
    \end{itemize}
}
\slide{Klassische Konzept-Theorie}{
    \begin{itemize}
        \item Struktur: Hierarchische Einordnung in Ober- und Unterbegriffe.
        \item Merkmale: Definition durch notwendige und hinreichende Merkmale.
        \begin{itemize}
            \item Notwendige Merkmale: Vom Oberbegriff vererbt (genus proximum).
            \item Hinreichende Merkmale: Spezifische Eigenschaften des Unterbegriffs (differentia specifica).
        \end{itemize}
        \item Vererbung: Unterbegriffe übernehmen automatisch alle Merkmale des Oberbegriffs.
    \end{itemize}
}
\slide{Semantische Netzwerke}{
    \begin{itemize}
        \item Aufbau: Graphische Darstellung von Wissen durch Knoten (Begriffe) und Kanten (Relationen).
        \item Types vs. Tokens:
        \begin{itemize}
            \item Types: Abstrakte Kategorien/Begriffe (z. B. "Mensch").
            \item Tokens: Spezifische Exemplare/Instanzen (z. B. "Max Mustermann").
        \end{itemize}
        \item Kantentypen (Relationen):
        \begin{itemize}
            \item isa (is a): Hierarchische Beziehung zwischen zwei Types (Unterbegriff ist Oberbegriff); Basis für Vererbung.
            \item instance-of: Beziehung zwischen Token und Type (Objekt gehört zu Kategorie).
            \item has-attribute / has-part: Zuordnung von Eigenschaften oder Teilen.
        \end{itemize}
        \item Probleme: Verifikationszeiten entsprechen nicht immer der Pfadlänge (ökonomisches Prinzip widerlegt); Unschärfe von Kategorien.
    \end{itemize}
}
\slide{Prototypen-Theorie (Rosch)}{
    \begin{itemize}
        \item Definition: Kategorien werden nicht durch feste Definitionen, sondern durch Ähnlichkeit zu einem Prototypen bestimmt.
        \item \b{Prototyp:} Der "typischste" Vertreter einer Kategorie; eine Idealisierung, die die meisten Eigenschaften vereint.
        \item \b{Typikalitätseffekt:} Manche Vertreter werden schneller als zugehörig erkannt als andere (z. B. Amsel vs. Pinguin als Vogel).
        \item Basiskategorien: Ebene mit höchstem Informationsgewinn; wird frühkindlich erworben (z. B. "Stuhl" statt "Möbel" oder "Bürostuhl").
        \item Kritik/Grenzen:
        \begin{itemize}
            \item Kontexteffekte: Kategoriengrenzen verschieben sich je nach Kontext (Tasse vs. Schüssel).
            \item Ad-hoc Kategorien: Spontan gebildete Kategorien für spezifische Ziele (z. B. "Dinge, die man aus einem brennenden Haus rettet"), die keinen festen Prototypen im Gedächtnis haben.
        \end{itemize}
    \end{itemize}
}
\slide{Komplexe Wissensrepräsentation}{
    \begin{itemize}
        \item \b{Schemata:} Komplexe, strukturierte Wissenseinheiten zur Interpretation von Umwelt-Erfahrungen.
        \item \b{Frames:} Rahmenstruktur zur Repräsentation von Objekten/Szenen; besteht aus Header (Name), Slots (Attribute mit Default-Werten) und Prozeduren . Ähnelt semantischen Netzwerken durch Vererbungshierarchien.
        \item \b{Scripts:} Stereotype Aktionsfolgen für soziale Situationen (z. B. Restaurantbesuch).
        \begin{itemize}
            \item Enthalten Rollen, Requisiten, Vor- und Nachbedingungen.
            \item Dienen der Verhaltenssteuerung, Entlastung (Automatisierung) und dem Schließen von Informationslücken (bridging inferences).
        \end{itemize}
    \end{itemize}
}